\chapter{Instrukcja uruchomienia projektu}

Projekt składa się z dwóch niezależnych gier: Nim oraz Reversi. Obie gry zostały napisane w języku Python i korzystają z biblioteki \texttt{tkinter} do obsługi interfejsu graficznego.

\section{Wymagania środowiskowe}

Do uruchomienia projektu wymagane jest środowisko Python w wersji 3.12 lub nowszej. Interfejs graficzny wykorzystuje bibliotekę \texttt{tkinter}, która jest częścią standardowej dystrybucji Pythona.

Wymagania:
\begin{itemize}
\item Python 3.12 lub nowszy,
\item terminal lub środowisko IDE, np. PyCharm.
\end{itemize}

\section{Struktura projektu}

Projekt został podzielony na dwa główne moduły: \texttt{nim\_project} oraz \texttt{reversi\_project}. Każdy z nich posiada podobną strukturę katalogów, co ułatwia analizę kodu i pokazuje spójny sposób organizacji projektu.

\begin{lstlisting}[caption={Uproszczona struktura projektu.}, label={lst:project-structure}]
SI_project/
|-- nim_project/
|   |-- main.py
|   |-- ai/
|   |-- gui/
|   |-- logic/
|   |-- models/
|
|-- reversi_project/
|   |-- main.py
|   |-- ai/
|   |-- gui/
|   |-- logic/
|   |-- models/
|
|-- .gitignore
\end{lstlisting}

Katalog \texttt{ai} zawiera logikę agentów decyzyjnych. Katalog \texttt{logic} odpowiada za zasady gry, poprawność ruchów i aktualizację stanu rozgrywki. W katalogu \texttt{models} znajdują się proste klasy danych, natomiast \texttt{gui} zawiera implementację interfejsu graficznego.

\newpage

\section{Uruchomienie gry Nim}

Aby uruchomić grę Nim, należy przejść w terminalu do głównego katalogu projektu \texttt{SI\_project}. Następnie należy wykonać polecenie:

\begin{center}
\texttt{python -m nim\_project.main}
\end{center}

Po uruchomieniu programu pojawia się ekran wyboru początkowego układu stert oraz trybu gry. Użytkownik może wybrać tryb Człowiek vs AI albo AI vs AI. W trybie Człowiek vs AI gracz sam wybiera żetony do zabrania, natomiast w trybie AI vs AI rozgrywka przebiega automatycznie.

\section{Uruchomienie gry Reversi}

Aby uruchomić grę Reversi, również należy znajdować się w katalogu głównym projektu \texttt{SI\_project}. Następnie należy wykonać polecenie:

\begin{center}
\texttt{python -m reversi\_project.main}
\end{center}

Po uruchomieniu aplikacji wyświetlana jest plansza Reversi oraz panel boczny. Użytkownik może wybrać tryb Człowiek vs AI albo AI vs AI. W trybie Człowiek vs AI dostępny jest wybór poziomu trudności oraz podpowiedzi. W trybie AI vs AI można ustawić osobną głębokość przeszukiwania dla AI 1 i AI 2.